\section{Knowing without eyes, and knowing without omniscience}
\label{sec:angel-knowledge}

\subsection{The knower, the power, and the act (Question 54)}

A musician, the musician's ability, and the performance are distinguishable even when they belong to one person. Aquinas asks an analogous question about an angel, while removing the bodily instrument. The angel exists, possesses a power of understanding, and actually understands. Those are not three substances, but neither are they three interchangeable names. The angel's understanding is an operation of a creature; it is not the creature's essence or its received act of existence. Only in God do being, understanding, and substance coincide.\footnote{\ST{I q.~54 aa.~1--3}. The musician comparison illustrates the distinctions; it does not supply an account of angelic habit acquisition.}

This protects the distinction established in question 50. Calling angels ``intellects'' does not make them self-subsisting acts of absolute understanding. It identifies the kind of life they possess. A created intellect can have real perfections beyond what constitutes its substance. Aquinas uses Dionysius's distinction of essence, power, and operation to articulate precisely this creaturely plurality.\footnote{\ST{I q.~54 a.~3}, corpus and ad~1--2; Dionysius, \emph{Celestial Hierarchy} XI, as received there.}

Article 4 removes another misleading comparison. Human intellect comes to understand material things through abstraction from sensory images. Aquinas distinguishes an active intellect that makes the intelligible content available and a possible intellect that receives it. An angel has no corresponding sensory workshop. Therefore it needs neither faculty in that specifically human sense. This does not mean that an angel cannot receive revelation or illuminate another angel. Those operations are different from abstracting intelligible forms out of images. A word such as ``active'' can be reused analogically without importing the whole human mechanism.\footnote{\ST{I q.~54 a.~4}, especially ad~2.}

Article 5 draws the bodily consequence. Angels do not possess imagination, sense memory, or sensation as organic faculties. When authorities speak of angelic experience or memory, Aquinas distinguishes the known content from the way it is known. An angel can know a particular occurrence without acquiring it through eyes, and retain knowledge without having a sensory memory. Such expressions do not turn angelic intellect into a human brain freed from its skull.\footnote{\ST{I q.~54 a.~5}.}

\angeldossiersettled{I q.~54, aa.~1--5}{Church teaching: angels possess intelligence (\emph{CCC} 330). Thomist position: distinctions of substance, power, and operation, and absence of the human abstraction apparatus.}{Dionysius, \emph{Celestial Hierarchy} XI and \emph{Divine Names} IV; Aristotle, \emph{De anima} III; Gregory, Gospel homily~29, as invoked in the articles.}

\subsection{How can knowledge arrive without sensation? (Question 55)}

Suppose an angel knows a tree. The tree does not enter the angel, and the angel does not become botanical matter. Aquinas calls the intelligible likeness by which something is known a \emph{species}. Here that word does not mean the metaphysical kind discussed in question 50. It means a form enabling knowledge. The same Latin term answers two different questions: what kind of being something is, and through what intelligible likeness it is known.

An angel cannot know every other thing merely by being itself. Its essence is finite; it does not contain every creature's perfection as the divine essence does. Article 1 therefore distinguishes the angel's own substance from the intelligible species through which other things are known. Article 2 identifies these species as connatural gifts from God, not impressions abstracted from bodily encounters. The Creator produces both things in their natural existence and intelligible likenesses proportioned to angelic minds.\footnote{\ST{I q.~55 aa.~1--2}.}

An innate species is thus not an imaginative photograph stored before birth. The \emph{De veritate} makes the alternatives explicit: the angel neither extracts its forms directly from material objects nor arbitrarily reshapes its essence to resemble whatever appears. The latter proposal would still need an intelligible principle guiding the transformation. Aquinas's account places that principle in the creature's original intellectual endowment.\footnote{Aquinas, \emph{De veritate} q.~8 a.~9, corpus.}

Article 3 asks why a superior angel knows through fewer and more universal species. In ordinary conversation, ``general knowledge'' can mean vague knowledge: someone knows there is a problem but cannot identify it. Aquinas means almost the opposite. A more universal intelligible principle can contain more distinct objects precisely understood. His comparison is with a strong intellect that grasps many implications from one explanation, where another requires each detail separately. Greater unity belongs to the \emph{medium of knowing}, not to an erasure of detail in the \emph{object known}.\footnote{\ST{I q.~55 a.~3}, especially ad~2--3.}

\angeldossier{I q.~55, aa.~1--3}{Thomist position: connatural intelligible species and their greater universality in higher angels.}{Dionysius, \emph{Divine Names} IV, VII and \emph{Celestial Hierarchy} XII; Augustine, \emph{Literal Meaning of Genesis} II.8, as received by Aquinas; \emph{De veritate} q.~8 a.~9.}{Aquinas rejects both knowledge of everything through the angel's own essence and acquisition of species from material things. Greater universality does not mean vaguer knowledge.}

\subsection{Self, another spirit, and God (Question 56)}

The fact that an angel needs species to know other things does not mean it needs a copy of itself to know itself. Its own immaterial substance is already intelligible and present to its intellect. Article 1 therefore attributes self-knowledge through its essence. This is knowledge of a finite created nature, not comprehension of every divine purpose concerning that nature.\footnote{\ST{I q.~56 a.~1}, corpus and ad~1.}

Knowledge of another angel is different. The other exists as an independent substance; its intelligible likeness exists within the knower. Aquinas calls this the difference between \emph{natural} and \emph{intentional} being. ``Intentional'' here does not mean imaginary or merely intended for the future: it names the presence of an object as known. The form through which one angel knows another is not a second actual angel living inside the first.\footnote{\ST{I q.~56 a.~2}, especially ad~3.}

There is a third case: knowledge of God. An angel's own nature bears an image of the Creator and therefore supplies natural knowledge of him. Yet no created likeness can disclose the divine essence as it is in itself. The natural knowledge is neither ignorance of God nor the beatific vision. Aquinas places it between human knowledge from sensible effects and supernatural face-to-face vision, while saying it is closer to knowledge through a mirror than to that vision. Even an unfallen angel cannot cross this difference by thinking harder.\footnote{\ST{I q.~56 a.~3}; compare \ST{I q.~12 aa.~4--5}.}

A striking textual moment occurs in article 1. An objection quotes Dionysius as saying that angels do \emph{not} know their own powers. Aquinas replies that a newer translation corrects the wording to an affirmative. He also explains the older wording in a restricted sense. The argument's authority depends on what the received text actually says, and a translated negative can change the apparent teaching.\footnote{\ST{I q.~56 a.~1 ad~1}, reporting Aquinas's comparison of translations of \emph{Celestial Hierarchy} VI.}

\angeldossier{I q.~56, aa.~1--3}{Thomist position: self-knowledge through essence and knowledge of other spirits through species. Common teaching: natural knowledge does not itself confer the vision of God.}{Augustine, \emph{Literal Meaning of Genesis} II; \emph{Book of Causes}; Romans~1:19--20; Dionysius and the translation distinction discussed in a.~1.}{An objection denies self-knowledge through a negative in an older Dionysian translation. Aquinas corrects the reading and distinguishes knowledge of powers from comprehension of their divine source.}

\subsection{A guardian must know this person (Question 57)}

An intellect that knew only ``humanity'' could not protect this traveller at this crossroads. Aquinas consequently joins the angel's knowledge of material natures to knowledge of singulars. The scriptural ministry is directed to actual persons, not abstract classes. The divine gift of intelligible species represents things with their individuating conditions; it does not merely duplicate human concepts abstracted from those conditions.\footnote{\ST{I q.~57 aa.~1--2}; Ps.~90:11 [Hebrew~91:11]; Heb.~1:14.}

That reach has three decisive limits. First, knowing what exists far away is different from knowing an undetermined future. Existing things can be represented as actual; future free events are not already present to an angel as they are to God's eternity. An angel can know necessary effects through their causes and form superior conjectures about probable effects. God can also reveal the future. None of this gives an angel natural possession of every future event. The \emph{De veritate} adds that a result contingent relative to one cause may be understood more certainly when a wider combination of causes is known. That still leaves genuinely free contingents distinct from determined natural effects.\footnote{\ST{I q.~57 a.~3}; \emph{De veritate} q.~8 a.~12, corpus.}

Second, another person's thought is not naturally transparent merely because an angel is immaterial. A bodily change or an outward act may disclose something about a person's concern; an angel can recognize such signs more acutely than we do. But a thought as freely held in the mind, and an affection as freely held in the will, are immediately open to God alone. Aquinas specifically distinguishes knowing the intelligible species another possesses from knowing the other's present use of them. Even the imaginative and sensitive life of a human being is not an infallible substitute for access to free thought.\footnote{\ST{I q.~57 a.~4}, corpus and ad~2--3; Jer.~17:9--10; 1~Cor.~2:11.}

Third, no angel can deduce God's free saving decisions from created nature. Mysteries of grace are known by revelation. Blessed angels see God, but seeing is not comprehending: they do not exhaust the divine wisdom. Accordingly they receive unequal and successive disclosures suited to their missions. Aquinas distinguishes their general knowledge of the Incarnation from knowledge of all its particular circumstances. A minister can truly know the purpose of a mission while receiving further instructions about its execution.\footnote{\ST{I q.~57 a.~5}, especially ad~1--2.}

\angeldossier{I q.~57, aa.~1--5}{Common teaching: angelic knowledge is finite, and secret hearts and future contingents are not naturally open to creatures. Thomist position: their knowledge of singulars through divinely given species and the detailed distinction of cognitive media.}{Psalm~90 [Hebrew~91]; Isaiah~41:23; Jeremiah~17:9--10; 1~Corinthians~2:10--11; Dionysius, \emph{Celestial Hierarchy} VII, as invoked in a.~5; \emph{De veritate} q.~8 a.~12.}{Aquinas rejects knowledge confined to universals, an exclusively causal knowledge of singulars, and the inference from immateriality to knowledge of every future event or secret thought.}

\subsection{Insight, succession, and the two lights of knowledge (Question 58)}

Human learning often requires a proof: first we understand the premises, then discover their conclusion. Aquinas does not picture angelic thought as the same process running extremely fast. In its natural field an angel sees the consequences within the principles. It knows a syllogism and its validity without acquiring a previously unknown conclusion by syllogizing. Likewise it knows what can be affirmed or denied of a nature without laboriously assembling subject and predicate. The distinction is between knowing a reasoning process and needing that process in order to know.\footnote{\ST{I q.~58 aa.~3--4}; \emph{De veritate} q.~8 a.~15, corpus.}

This simplicity does not mean that the angel considers everything at once. Article 1 distinguishes not yet possessing knowledge from possessing it without attending to it now. Natural species are already given; their exercise can be successive. Revelation can also disclose what was previously unknown. Article 2 adds that objects grasped under one intelligible species can be known together; those requiring distinct species cannot be known together in that way. The simultaneous vision of things in the Word extends to what God actually shows, not automatically to everything God knows.\footnote{\ST{I q.~58 aa.~1--2}; read with \ST{I q.~57 a.~5}; \emph{De veritate} q.~8 a.~14, corpus.}

Can such an intellect be wrong? Aquinas's answer distinguishes natural understanding from presumptuous judgment about what exceeds it. The angel does not misread a nature within its proper intellectual field. A fallen angel may nevertheless judge what God will do as though natural conditions exhausted the possibilities: that a dead person will not rise, for example. The defect is no shortage of cleverness. It is the disordered extension of natural knowledge beyond its warrant. Good angels submit their judgments to divine wisdom.\footnote{\ST{I q.~58 a.~5}; \emph{De malo} q.~16 a.~6, corpus.}

The last two articles receive Augustine's ``morning'' and ``evening'' knowledge. Morning names knowing creatures in their source, the Word; evening names knowing their created existence. These are not working shifts or alternating periods of blindness. A creature considered in its own nature remains good and intelligible, but its light is less than the uncreated source. The good angel refers that created knowledge back to praise: evening leads to morning rather than closing into night.\footnote{\ST{I q.~58 a.~6}, especially ad~1--3.}

Article 7 asks whether two kinds of knowledge have actually been identified. If the angel knows a creature's own existence \emph{through the Word}, morning and evening differ by the aspect known within one medium. If it knows that existence \emph{through connatural species}, the media differ. Aquinas favors the latter as Augustine's meaning. Their coexistence is possible because the lower operation is ordered to the higher. Glory does not require the annihilation of created understanding.\footnote{\ST{I q.~58 a.~7}.}

\angeldossier{I q.~58, aa.~1--7}{Thomist position: nondiscursive knowing, its limits, and the account of error. Disputed theological interpretation: Augustine's morning--evening reading and the relation of its cognitive media.}{Dionysius, \emph{Divine Names} VII; Augustine, \emph{Literal Meaning of Genesis} IV.22--24, 31--32, as received in the articles; Aquinas, \emph{De veritate} q.~8 aa.~14--15 and \emph{De malo} q.~16 a.~6.}{Aquinas rejects both human-style discursion and an unrestricted simultaneity of natural knowledge; in a.~7 he distinguishes two defensible senses of ``evening knowledge.''}
